\documentclass[12pt,a4paper]{article}

% Encoding
\usepackage[utf8]{inputenc}
\usepackage[T1]{fontenc}

% Layout
\usepackage[margin=2.5cm]{geometry}
\usepackage{setspace}
\onehalfspacing

% Packages
\usepackage{graphicx}
\usepackage{amsmath}
\usepackage{hyperref}
\usepackage{parskip}
\usepackage{titlesec}
\usepackage{float}      % <-- ÖNEMLİ (figure sabitleme için)
\usepackage{caption}

% Section style
\titleformat{\section}{\large\bfseries}{\thesection.}{0.5em}{}
\titleformat{\subsection}{\normalsize\bfseries}{\thesubsection.}{0.5em}{}

\begin{document}

% ─── COVER PAGE ─────────────────────────────
\begin{titlepage}
    \centering

    \vspace*{3cm}

    {\Large \textbf{CSE344 -- Systems Programming}}\\[1.5cm]

    {\LARGE \textbf{Homework \#2}}\\[0.5cm]
    {\large procSearch Report}\\[2cm]

    \begin{tabular}{rl}
        \textbf{Name:} & Ayşe Feyza SERBEST \\
        \textbf{Student ID:} & 220104004052 \\
        \textbf{Course:} & CSE344 \\
        \textbf{Homework:} & \#2 \\
        \textbf{Date:} & \today \\
    \end{tabular}

    \vfill
\end{titlepage}

\tableofcontents
\newpage

% ─── INTRODUCTION ─────────────────────────────
\section{Introduction}

This program is a multi-process file search utility written in C. 
It searches files inside a directory using multiple worker processes 
and reports matching files based on a given pattern and optional size filter.

The main goal of the design is to efficiently utilize multiple processes 
while maintaining correct inter-process communication and safe signal handling.

% ─── DESIGN DECISIONS ─────────────────────────
\section{Design Decisions}

\subsection{Directory Partitioning}

Directories are divided among worker processes using a round-robin approach. 
This ensures that the workload is distributed evenly regardless of directory order.

\subsection{Signal Handling Design}

Signal handling was one of the most challenging parts of this project.

Handling signals asynchronously requires careful design because most standard 
functions are not safe to use inside signal handlers.

To solve this problem:

\begin{itemize}
    \item Global variables were used to share state between the main program and signal handlers.
    \item Only async-signal-safe functions such as \texttt{write()} were used.
    \item Workers respond to \texttt{SIGTERM} by sending partial results before exiting.
\end{itemize}

Using global variables significantly simplified communication between the signal 
handler and the main execution flow.

\subsection{Avoiding Async-Signal-Safety Issues}

To avoid undefined behavior:

\begin{itemize}
    \item No \texttt{printf()} was used inside signal handlers.
    \item Only \texttt{write()} system call was used.
    \item Shared variables were defined as \texttt{volatile sig\_atomic\_t}.
\end{itemize}

\subsection{Regex Implementation}

The regex functionality was another challenging part.

Instead of using a full regex library, a simplified custom matching logic was implemented.

The main difficulty was handling repetition patterns like \texttt{c+}, 
which required careful recursive matching.

% ─── SCREENSHOTS ─────────────────────────────
\section{Screenshots}

\subsection{Test Case 1}
\begin{figure}[H]
    \centering
    \includegraphics[width=0.8\textwidth]{a.png}
    \caption{Regex matching}
\end{figure}

\subsection{Test Case 2}
\begin{figure}[H]
    \centering
    \includegraphics[width=0.8\textwidth]{b.png}
    \caption{15 bytes}
\end{figure}

\subsection{Test Case 3}
\begin{figure}[H]
    \centering
    \includegraphics[width=0.8\textwidth]{c.png}
    \caption{20 bytes}
\end{figure}

\subsection{Test Case 4}
\begin{figure}[H]
    \centering
    \includegraphics[width=0.8\textwidth]{d.png}
    \caption{No matching files found.}
\end{figure}

\subsection{Test Case 5}
\begin{figure}[H]
    \centering
    \includegraphics[width=0.8\textwidth]{e.png}
    \caption{Avoid zombie processes.}
\end{figure}

\subsection{Test Case 6}
\begin{figure}[H]
    \centering
    \includegraphics[width=0.8\textwidth]{f.png}
    \caption{monitor.log}
\end{figure}

% ─── CONCLUSION ─────────────────────────────
\section{Conclusion}

This project helped deepen my understanding of:

\begin{itemize}
    \item Process management with \texttt{fork()}
    \item Inter-process communication using pipes
    \item Signal handling in Unix systems
\end{itemize}

The most difficult parts were signal handling and regex implementation.

Signal handling was particularly challenging due to async-signal-safety constraints. 
Using global variables helped manage shared state safely.

Regex implementation required careful thinking to correctly match patterns 
with repetition.

Despite these challenges, the project provided valuable experience 
in system-level programming.

\end{document}